\chapter{知识的迁移与适应}
\label{chap:ket-transfer-adaptation}

一个已经训练好的花卉分类器，要怎样帮助另一个植物园完成新的工作？新园区可能要求识别来源数据中没有的物种，也可能要求判断叶片病害，或在照片中标出每朵花的位置。原模型提取的颜色、纹理和形状信息仍可能有用，但把它原样运行一次，未必能得到新任务需要的答案。我们需要找出来源经验与目标问题之间的联系，再确定目标应补充哪些证据、允许怎样计算，才能形成可验证的预测能力。

第~\ref{chap:ket-update-maintenance}章已经讨论怎样改变系统并维持仍需提供的能力；本章把重点放在来源知识怎样帮助一个新的学习目标。学习新物种与保留旧物种可以是同一次操作的两项责任，却要分别验证。沿用前文的损失、风险和参数更新概念，我们先确定目标语义、信息权限及历史任务的联系，再考察能为适应准备什么，最后区分目标端构建参照、读取上下文和优化参数的操作。这不是每个方案都要重新训练一遍的流程：已有模型足够时，可以直接进入目标使用。

目标只提供几张标注图像，与准备阶段利用许多历史任务，是不同层面的条件。少样本和零样本说明目标缺少哪类学习材料；\term{元学习}（meta-learning）则利用历史任务训练学习过程，使共享知识在接受新任务证据后能够用于预测。它可以学习表示、起点或证据读取方式，不限定目标端是否做梯度更新。因而同一方案可以既是少样本学习，又采用元学习准备。

\section{可迁移的目标任务}
\label{sec:ket-transfer-targets}

新园区的要求必须先变成可评价的预测问题。识别物种时，一次错误可以是选错类别；定位花朵时，还要评价位置与对象匹配，二者不能共用一个没有说明含义的分数。设来源的输入空间、真实答案空间和联合分布为$X_S,Y_S,\mathcal P_S$，目标对应为$X_T,Y_T,\mathcal P_T$。对目标预测器$f$，用$\ell_T$给每个预测与真实答案之间的差异计分，再在未来目标样本上取平均，得到
\begin{equation}
  R_T(f)=\mathbb E_{(\bm{x},y)\sim\mathcal P_T}
  \bigl[\ell_T(f(\bm{x}),y)\bigr].
  \label{eq:ket-transfer-target-risk}
\end{equation}
其中$(\bm x,y)\in X_T\times Y_T$表示目标输入与真实答案，$f(\bm x)$可以是类别、分数、概率或结构化输出，只须与$\ell_T$的预测接口相容；假定损失的期望存在。期望采用$\mathcal P_T$，是因为要判断未来目标样本上的平均错误，而非重复衡量来源训练集。固定任务、损失、允许信息和预算后，较小的$R_T$才表示较少的目标错误；它本身还没有说明收益来自哪项来源知识。输入分布、输出语义或评价标准改变时，都要重新核对这种联系。

\subsection{预测目标与输出语义的确定}
\label{sec:ket-transfer-output-semantics}

先比较三种仍以花卉照片为输入的目标。新增物种以后，若一张图仍只输出一个物种名，需要扩展类别集合并学习新的分类边界。病害判断改的则是标签含义：同一物种的不同植株可以有不同病害状态，原来的物种标签没有给出这种区分。把输出编号换成病害名，并不能补上缺失的监督与视觉信息。

定位还把预测单位从整幅图像变成图中的对象。输出须表达花朵数量、各自位置和类别，训练时须说明哪个预测对应哪个标注对象。只增加向量长度，没有规定这些坐标的意义、变长对象的处理和匹配损失，便没有定义检测任务。分割以像素区域为单位，生成以输出序列为单位，也各自需要相应的输出结构和监督；来源表示能够复用到哪里，要按这些新的预测单位检查。

回到单幅图像分类，类别编号可比较以前，类别含义必须先对齐。令$C_S,C_T$为源、目标实际可能出现的语义类别，$C_{\cap}=C_S\cap C_T$为共享类。来源按物种、目标按园艺品种划分时，同名类别可能范围不同，需要先明确映射。表~\ref{tab:ket-transfer-class-relations}再固定一种训练协议：有带标签源样本和无标签目标样本，目标标签只用于方案锁定后的评价。此时类别集合的关系决定哪些源类应迁移、哪些目标输入需要拒识。监督权限是另一个维度；有少量目标标签或只能读取源模型而不能访问源数据时，须单列这些条件。

\begin{table*}[htbp]
  \centering
  \small
  \setlength{\tabcolsep}{3pt}
  \begin{tabularx}{\linewidth}{p{19mm}p{25mm}p{29mm}XX}
    \toprule
    设定 & 类别关系 & 代表性训练访问 & 目标端判断 & 额外风险 \\
    \midrule
    完全共享 & $C_T=C_S$ & 带标签源样本；无标签目标样本 & 在共同类别中分类 & 拍摄条件变化仍可破坏原分类器 \\
    部分领域自适应 & $C_T\subsetneq C_S$ & 同左；训练时不揭示目标标签 & 迁移目标实际具有的类别 & 源中特有类别可能被错误地对齐到目标 \\
    开放集领域自适应 & 典型设定为$C_S\subsetneq C_T$ & 同左；目标私有类不提供类别名 & 共享类分类，目标私有类拒识为未知 & 未知不能被强迫匹配到已知类 \\
    通用领域自适应 & 两侧均可有私有类 & 同左；共享范围预先未知 & 共享类分类，其余目标样本输出未知 & 共享范围与未知比例都需要估计 \\
    \bottomrule
  \end{tabularx}
  \caption{同一图像分类协议下的类别关系与预测责任。四行均允许带标签源样本和无标签目标样本；集合关系在语义对齐后判断。开放集一行采用$C_S\subsetneq C_T$的常用设定，通用设定不预定包含关系。输出“未知”只承认样本不属于共享已知类，不给出各目标私有类的名称。}
  \label{tab:ket-transfer-class-relations}
\end{table*}

目标只包含部分源类时，不能把每个源类都当成应在目标中找到的对象。\term{部分领域自适应}（partial domain adaptation）在这种包含关系下减少源私有类对迁移的干扰。部分对抗领域自适应（Partial Adversarial Domain Adaptation, PADA）从目标预测中估计源类别权重，再按权重调整源分类损失与领域对抗损失。领域对抗训练试图让特征难以被区分为来自源还是目标；加权则限制哪些源内容参与这种对齐。权重仍依赖当前预测，若共享类被早期模型漏判，其贡献也可能被错误压低\scite{ket-cao2018partial}

目标含有源中没有的物种时，问题反过来发生在目标侧：把所有目标特征拉向源类，会迫使未知物种接受一个已知名。\term{开放集领域自适应}（open set domain adaptation）允许输出“未知”。Saito等让分类器与特征生成器在目标输入上形成对抗，使一部分目标样本与已知类对齐，另一部分与已知类分离；这种机制需要检验未知拒识，而不能只评价共享类准确率\scite{ket-saito2018openset}

若连哪些类共享也事先不知道，\term{通用领域自适应}（universal domain adaptation）不再规定两侧标签集的包含关系。它需要根据模型可获得的证据判断哪些样本属于共享内容，并把其余目标样本标为未知。这比已知包含关系多了一项共享范围判断；“通用”放宽的是类别集合关系，不是对任意新任务给出成功保证\scite{ket-you2019universal}

这里还须分清三个结果。\term{未知拒识}（unknown rejection）仅表示现有类别目录不能为这张图提供合适答案。把拒识图像聚成几个群组，是进一步发现输入结构，但群组可能按背景而非物种形成。建立新类别能力还需要确定群组的语义、补充代表性区分证据，并在新样本上检验分类规则。例如可以先聚类，再由专家命名与补标，最后训练分类器；一个统一的“未知”输出没有完成这条后续过程。

以来源$\{\text{菊花},\text{月季},\text{鸢尾}\}$、目标$\{\text{月季},\text{鸢尾},\text{新种甲}\}$为例，共享内容是月季与鸢尾，菊花不应成为目标必须选择的答案。目标只给无标签照片而没有新种甲的描述或标注时，能够争取的是拒识新种甲，不能凭空给它建立有语义的分类器。后续描述与标注才为这一新增能力提供依据。排除菊花对目标的干扰，也不等于要求系统遗忘菊花；旧园区仍需识别它时，应另按第~\ref{sec:ket-update-retention-validation}节验证保留。

这些语义变化可以与分布变化同时存在。类别不变而相机或季节改变，主要改变的是观测分布；物种标签改为病害标签，即便输入分布不变，预测任务也已经不同。输入对齐不能替代新的标签联系。即使输入和标签空间固定，相同的可见数据也可能对应相反的目标标签规律，相应不可识别反例可参考\BookChapterRef{chap:transfer-learning-theory}{迁移学习和持续学习的可学习性}。新任务也须先提供连接来源与目标的证据，再使用第~\ref{sec:ket-update-weighting}节的加权或式~\eqref{eq:ket-update-crossmodal}的配对约束，并重新检查其覆盖与信息保留条件。本书将常被称为“领域适应”的相关问题统一称为领域自适应。

\subsection{目标信息与适应权限的确定}
\label{sec:ket-transfer-evidence}

对新种甲和新种乙，植物园提供什么，决定模型能从哪里获得区分依据。文字描述说明类别应具有什么特征，带标签图像把类别含义与实际观测连接起来，无标签当地照片则提供输入频率与局部结构。只给描述、每类给几张标注图像、再开放一批无标签照片，是三个信息条件。执行时都只有一次前向预测，也不使它们成为相同证据下的比较。

\term{少样本学习}（few-shot learning）描述目标只有少量学习用标注的条件。分类任务常规定$N$类、每类$K$例，写作$N$-way $K$-shot，总支持样本数为$NK$。这个写法没有交代查询数量、类间不平衡、图像重复和增强方式，比较时仍须补全。若只有描述而没有目标类训练图像，通常称为\term{零样本学习}（zero-shot learning），但“零”必须绑定缺失对象：没有新类图像仍可有名称和属性，没有语言任务示范仍可有任务指令，也可能在预训练中接触过相关内容。不同零样本协议不能仅按名称比较信息需求。

为把这些条件写清，可以约定目标证据包由类别描述$D_T$、带标签支持集$S_T$、许可的无标签集$U_T$与已到达的交互历史$H_T$组成。某一项可以为空，方案也不必用尽所有许可材料。比较时固定这个包及其可见时机，再记录来源预训练数据、专家整理描述的投入、检索库中的示范和额外选模标签。否则，较低目标风险可能来自更多知识，而不是适应操作本身。

证据数量相同以后，还要约定怎样使用查询输入。\term{归纳式预测}（inductive prediction）在适应结束后逐个处理新输入，不让其他查询改变当前答案。\term{直推式预测}（亦称转导式预测，transductive prediction）则允许联合利用整批无标签查询，例如按其聚类结构调整边界。后一协议增加了可用的输入关系，但没有开放隐藏标签；这项信息增益不能全记在优化器名下。部署时只有单张照片到达，依赖整批查询的计算也未必可原样执行。

参数读写、状态保留和反馈时机同样属于协议。模型能否更新骨干，能否跨请求保存记忆，能否先向专家查询，都改变了当前可执行的学习过程。主动查询怎样获取并校验新信息，由第~\ref{part:efficient-knowledge-use}部分讨论；这里记录查询次数、费用和反馈到达时刻。后来得到的正确标签可以用于后续预测，不能倒用于已经发生的无反馈预测，再把更新结果计作原协议的测试成绩。

在这些边界内，可以比较给定标注数$n$、更新步数$T$与计算预算$B$时的目标风险，或达到预定风险阈值所需的投入。风险越低越好，成本则须使用一致的单位；匹配的普通目标训练是判断迁移收益的对照。描述整理、示范检索、提示选择和验证调参都应进入账目。只数反向传播次数，会漏掉不更新权重的方案仍然承担的主要适应费用。

\subsection{历史任务与目标任务关系的确定}
\label{sec:ket-transfer-task-splits}

为新地区学习做准备，需要的不只是许多图像，还要有足以反映未来学习问题的历史任务。一个任务$\mathcal T_k$规定输入输出语义、联合分布$\mathcal P_k$、损失和信息权限。记$\Pi$为这些任务的抽样分布：先抽取“在某地区区分哪几种花”的问题，再从该问题中抽取图像，这是任务与样本的两层抽样。历史任务若共同依赖花瓣形状，可能支持相应表示的复用；若只共享背景，未来换地区时这种联系可能失效。准备方案须明确它预期保留哪种规律。

任务内也需要分离学习与检查。用于适应的\term{支持集}（support set）$S_k$给出当前任务的输入答案对；\term{查询集}（query set）$Q_k$检验据此形成的规则能否用于其他样本。二者可含相同类别，但须按实际评价单位隔离，不能将同一原图的增强副本分别放入两组，植株级评价还须隔离同一植株记录。任务间再分为元训练、元验证与最终测试任务；若声称迁移到未见类，三组类别也须分离。

要比较哪种共享准备更利于未来适应，不能只计算它在历史输入上直接预测的损失。把共享准备记为$\bm\phi$，支持证据形成任务预测器的过程记为$A_{\bm\phi}(S_k)$，用非空$Q_k$上的平均损失$\hat R_{Q_k}$评价该预测器。对任务抽样及任务内抽样取平均，得到一个典型准备目标
\begin{equation}
  \min_{\bm{\phi}}
  \mathbb E_{\mathcal T_k\sim\Pi}
  \mathbb E_{S_k,Q_k\mid\mathcal T_k}
  \bigl[\hat R_{Q_k}(A_{\bm{\phi}}(S_k))\bigr],
  \label{eq:ket-transfer-meta-objective}
\end{equation}
最小化在预先选定的共享参数范围内进行，期望存在时，它比较的是任务经支持适应后的平均查询错误。有限历史任务以采样平均近似这两层期望，不保证精确找到最小值。元训练查询标签参与$\bm\phi$的更新，因此只在任务内承担检查作用，并不是最终独立测试。模型结构、步长与停止条件由元验证任务选择；留出目标的查询答案直到方案冻结后才用于评分。图~\ref{fig:ket-transfer-data-boundaries}中历史虚线反馈存在、目标评分反馈不存在，正是两层数据边界的区别。

\begin{figure*}[htbp]
  \centering
  % Editable schematic, designed at 166 mm width; no measured data.
% Solid arrows: evidence use / prediction. Dashed arrows: training feedback.
\begin{tikzpicture}[
  x=1mm,y=1mm,
  font=\figurefont\footnotesize,
  text=BookInk,
  >={Latex[length=2mm,width=1.4mm]},
  evidence/.style={draw=BookBlue,fill=FigureInput!10,line width=0.8pt,
    text width=24mm,minimum height=12mm,align=center,inner sep=2mm},
  model/.style={draw=BookTeal,fill=FigureModel!10,line width=0.8pt,
    text width=26mm,minimum height=12mm,align=center,inner sep=2mm},
  score/.style={draw=BookAmber,fill=FigureLoss!10,line width=0.8pt,
    text width=28mm,minimum height=12mm,align=center,inner sep=2mm},
  flow/.style={draw=BookInk,line width=0.9pt,-{Latex},rounded corners=1mm},
  feedback/.style={draw=BookAmber,line width=0.9pt,dashed,-{Latex},
    rounded corners=1mm}
]
  \node[anchor=west,font=\figurefont\small] at (2,15) {历史任务：允许跨任务训练反馈};
  \node[evidence] (support-one) at (16,0) {$S_1$：月季、鸢尾\\支持图像与标签};
  \node[model] (adapt-one) at (59,0) {用共享准备$\bm\phi$\\形成任务预测器1};
  \node[score] (query-one) at (101,0) {$Q_1$：其他图像\\查询预测与损失};

  \node[evidence] (support-two) at (16,-29) {$S_2$：菊花、百合\\支持图像与标签};
  \node[model] (adapt-two) at (59,-29) {用共享准备$\bm\phi$\\形成任务预测器2};
  \node[score] (query-two) at (101,-29) {$Q_2$：其他图像\\查询预测与损失};
  \node[model,text width=25mm,minimum height=19mm] (prepare) at (150,-14.5)
    {跨任务更新\\共享准备$\bm\phi$};

  \draw[flow] (support-one.east) -- (adapt-one.west);
  \draw[flow] (adapt-one.east) -- (query-one.west);
  \draw[flow] (support-two.east) -- (adapt-two.west);
  \draw[flow] (adapt-two.east) -- (query-two.west);
  \draw[feedback] (query-one.east) -- (123,0) |- (prepare.west);
  \draw[feedback] (query-two.east) -- (123,-29) |- (prepare.west);
  \node[anchor=west,text=BookMuted] at (2,-43)
    {任务1、2均从同一轮准备出发；更新后进入下一轮历史任务};

  \draw[draw=BookRule,line width=0.8pt] (2,-50) -- (165,-50);
  \node[anchor=west,font=\figurefont\small] at (2,-59)
    {留出目标：方案冻结后，只用支持证据适应};

  \node[evidence] (support-target) at (16,-76)
    {$S_T$：新种甲、乙\\允许的图像与标签};
  \node[model] (adapt-target) at (59,-76)
    {目标任务预测器\\由$S_T$形成};
  \node[model,text width=23mm,fill=FigureOutput!8] (predict-target) at (103,-76)
    {$Q_T$预测输入\\目标预测};
  \node[score,text width=24mm] (score-target) at (151,-76)
    {$Q_T$隐藏答案\\最终评分};
  \draw[flow] (support-target.east) -- (adapt-target.west);
  \draw[flow] (adapt-target.east) -- (predict-target.west);
  \draw[flow] (predict-target.east) -- (score-target.west);
  % Frozen transfer is a one-way route outside the historical task lanes.
  \draw[flow] (prepare.south) -- (150,-47) -- (167,-47)
    -- (167,-64) -- (59,-64) -- (adapt-target.north);
  \node[anchor=east,fill=white,inner sep=1mm,text=BookTeal]
    at (148,-64) {冻结准备结果};
  \node[anchor=west,text=BookMuted] at (2,-92)
    {目标评分没有返回适应或模型选择的箭头};
\end{tikzpicture}

  \caption{共享准备如何接收历史反馈而隔离目标答案。实线表示支持证据、冻结准备和预测的使用，虚线表示历史查询损失更新共享准备。历史任务各自从同一轮准备出发；留出目标只用许可的支持证据形成预测器，查询答案只进入最终评分。图中的花卉类别为教学划分，目标还可包含独立留出的领域变化。}
  \label{fig:ket-transfer-data-boundaries}
\end{figure*}

Baxter用归纳偏置学习分析这种跨任务选择。候选假设空间规定每个任务可以拟合哪些函数，任务内数据帮助选函数，多个任务的经验则帮助选空间。增加一项任务的图像，可以改善该任务的估计，却没有增加独立任务关系的证据；增加独立任务，才有助于判断所选限制是否适用于未来任务。这种论证依赖任务联系以及候选空间族的容量控制，否则共享准备仍可能只记住有限历史任务\scite{ket-baxter2000-bias}

历史覆盖的检验还要区分类别与领域。相同设备上换一组物种，只能检验未见类；从温室照片换到野外红外影像，还改变了观测与领域。Meta-Dataset纳入多种数据来源，并改变任务类别数与支持数量，使评价不局限于固定五分类；要声称跨领域迁移，仍须将相应目标领域整体留出\scite{ket-triantafillou2019metadataset}成绩变化应分别对应类别新颖性、领域变化、标注数量和类别区分难度，不能把多个变量合成一个未定义的“任务更难”。

历史任务有限时，可以尝试合成任务来扩展准备材料。任务插值混合历史任务的中间特征；共享标签语义时可同步插值标签，不共享时需要配对类别并重新规定任务内标签。支持与查询分别构造，查询答案不能进入支持输入。插值定义了新的训练问题，却没有证明它包含真实新地区或新物种的变化；收益仍由独立目标任务检验\scite{ket-yao2021interpolation}

即便查询成绩很高，也要检查模型是否实际读取了支持信息。如果月季在所有历史任务中都标为1，模型可能直接记住这种固定映射，无须从本次支持集学习。为检查这条捷径，可以在每个任务内对支持集与查询集使用同一正确标签映射，同时在不同任务之间独立重新分配类别标签；替换支持集或移除输入暴露的任务身份，则分别检查另外两条记忆路径。Yin等用信息约束研究如何限制参数或查询输入中的任务记忆，促使预测依赖支持信息；这种处理不保证消除所有任务中的记忆路径\scite{ket-yin2019memorization}数据隔离的统一规则与任务层重复分析分别见第~\ref{sec:ket-evaluation-splits}、\ref{sec:ket-evaluation-repeats}节。

\section{可迁移知识的准备}
\label{sec:ket-transfer-preparation}

目标与历史联系确定后，需要问已有模型还缺哪种可复用计算。少量花卉图像需要一个能比较它们的表示；只有类别描述时，需要把语义与视觉连接起来；允许优化时，需要合适的起点、先验或更新规则；直接读示范时，则需要证据到预测的条件映射。下面按这些功能检查准备内容，它们可以共存，不是五类互斥模型。已有预训练表示足够时无需另做元训练；只有未来任务的联系和实际使用数量支持投入时，专门学习适应过程才有理由。

\subsection{可迁移表示的准备}
\label{sec:ket-transfer-representation}

每类只有几张图像时，首先检查来源编码器是否保留了区分类别所需的线索。记$\bm z=h_{\bm\theta}(\bm x)\in\mathbb R^d$为图像表示，$d$为特征维数。若目标类别在该表示中可分，标注主要用来确定类别区域；若两个不同病害输入被编码成同一向量，任何只读取这个向量的任务头都无法区分它们。要分开检查表示与读取规则，须固定来源数据、增强和目标划分。在同一冻结表示上比较线性分类、近邻和类别均值，只能说明这些读取规则在该表示上的结果；要比较表示，则须固定读取规则及训练预算，只改变来源编码器。表示与读取规则仍可能交互，因此任一对照只支持当前固定接口下的结论。

普通监督预训练是这项比较不能省略的来源。Chen等在一致实现中比较少样本分类方法，并检查普通微调的跨领域基线；Tian等考察冻结表示后拟合线性分类器及自蒸馏的影响。这些研究提醒我们，将基准差异归于元学习之前，应控制表示和实现条件，不能反过来推出所有目标都只需固定特征。自蒸馏还多了一次教师学生训练，其机制见第~\ref{sec:ket-integration-student-training}节，费用也须计入准备\scite{ket-chen2019closer,ket-tian2020embedding}

若准备目标就是依靠少量支持图像分类，可以训练表示，使同类支持点能够作为查询的参照。\term{原型网络}（Prototypical Networks）把每类支持表示的均值作为代表点，再按查询到代表点的距离分配类别概率。对任务$k$的$N$个候选类，令非空$S_{k,c}$为类别$c$的支持子集。采用平方欧氏距离时，代表点与概率写为
\begin{align}
  \bm c_{k,c}
  &=\frac{1}{|S_{k,c}|}
    \sum_{(\bm x_i,y_i)\in S_{k,c}}h_{\bm{\theta}}(\bm x_i),
    \label{eq:ket-transfer-prototype}\\
  p_{\bm{\theta}}(y=c\mid\bm x,S_k)
  &=\frac{\exp\{-\|h_{\bm{\theta}}(\bm x)-\bm c_{k,c}\|_2^2/\tau\}}
    {\sum_{c'=1}^{N}\exp\{-\|h_{\bm{\theta}}(\bm x)-\bm c_{k,c'}\|_2^2/\tau\}}.
    \label{eq:ket-transfer-prototype-probability}
\end{align}
其中$\bm c_{k,c}\in\mathbb R^d$是类别原型，$\tau>0$为温度，标准形式可取1。负距离越大，表示查询越近；指数归一化将$N$个候选的分数转成和为一的概率。选择均值也有具体依据：$\sum_i\|\bm z_i-\bm c\|_2^2$对$\bm c$的导数为$2|S_{k,c}|\bm c-2\sum_i\bm z_i$，置零给出上述均值，且二次项保证这是最小点。它回答的是怎样以单点减少类内平方偏离，不保证多峰类别可由一个点充分代表\scite{ket-snell2017-prototypical}

准备时反复模拟这种少样本使用。一次任务从历史类别抽$N$类，每类抽$K$个支持样本，再从隔离的剩余图像抽取查询。当前编码器同时处理两组输入，支持形成原型，查询按式~\eqref{eq:ket-transfer-prototype-probability}预测。要提高正确查询类别的概率，以非空查询集上的平均负对数概率为训练目标：
\begin{equation}
  L_k(\bm{\theta})
  =-\frac{1}{|Q_k|}\sum_{(\bm x,y)\in Q_k}
  \log p_{\bm{\theta}}(y\mid\bm x,S_k).
  \label{eq:ket-transfer-episodic-loss}
\end{equation}
梯度通过查询表示和支持均值共同更新$\bm\theta$，让表示适合按支持原型分类。编码器一变，原型就须以当前支持表示重算，不能始终使用准备开始时的均值。历史任务不断重采后，最终保存的是编码器及距离、温度等约定；未来类别的原型还要由未来支持集建立。

单个均值可能抹掉同类图像中的不同模式。\term{匹配网络}（Matching Networks）改为保留逐个支持参照，让相似样本的标签贡献更大。令$f_{\bm\theta}(\bm x,S_k)$与$g_{\bm\theta}(\bm x_i,S_k)$为维度相容的查询和支持表示，$\operatorname{sim}$为规定的相似度，值越大越相似。先把相似度归一化成支持权重，再对同标签权重求和：
\begin{align}
  s_i(\bm x,S_k)
  &=\operatorname{sim}\bigl(f_{\bm{\theta}}(\bm x,S_k),
       g_{\bm{\theta}}(\bm x_i,S_k)\bigr), \nonumber\\
  a_i(\bm x,S_k)
  &=\frac{\exp\{s_i(\bm x,S_k)/\tau\}}
  {\sum_{j=1}^{|S_k|}\exp\{s_j(\bm x,S_k)/\tau\}}, \nonumber\\
  p_{\bm{\theta}}(y=c\mid\bm x,S_k)
  &=\sum_{i=1}^{|S_k|}a_i(\bm x,S_k)\,\mathbf 1\{y_i=c\}.
  \label{eq:ket-transfer-matching}
\end{align}
这里$i$遍历支持记录，$\mathbf1\{y_i=c\}$只将类别$c$的权重留下，温度$\tau>0$控制权重集中程度。原论文用余弦相似度，还研究让编码依赖整个支持上下文：双向循环网络读取支持排列，查询通过多次注意力读取支持表示。同一图像的编码因此可随当前支持内容变化，也增加排列处理与循环计算，不能按独立固定编码器核算\scite{ket-vinyals2016matching}

将式~\eqref{eq:ket-transfer-matching}代入同一查询损失，即可按历史任务训练匹配网络。固定独立表示时，逐例匹配的参照数量随支持样本数增长，原型则只保留每类一个点，减少存储与比较次数，但改变了类别形状假设。每类一例且编码器、比较分数相同时，均值就是该例，两种规则一致；多例时若支持结构特殊，两者也可能对某次查询给出相同结果，不能据此宣称普遍等价。余弦与平方距离不同，或只有一方用上下文编码时，同样不能仅据样本数宣称等价。第~\ref{sec:ket-transfer-references}节进一步计算目标参照，并检查多例情形的差异。

\subsection{可迁移语义联系的准备}
\label{sec:ket-transfer-semantic-preparation}

没有新物种的训练图像，就不能为它直接求支持均值。如果目标改为提供“花瓣细长、叶缘有锯齿”的描述，还需要一个从图像辨认这些特征的计算过程。类别描述回答应寻找什么，来源监督回答怎样从观测中识别它；两者连接以后，描述才能成为新类预测的依据，而不是仅给输出换一个名称。

先用$m$个二值属性说明这种联系。类别$c$的描述为$\bm a_c\in\{0,1\}^m$，第$j$项说明属性是否存在。\term{直接属性预测}（direct attribute prediction, DAP）用来源属性监督训练$q_j(\bm x)\approx p(a_j=1\mid\bm x)$，再用预测属性匹配目标描述。保存的是$m$个属性预测器，输出维数不随目标类别数增加。来源监督可以是逐图属性标注，也可以把类别描述赋给该类每张图；后一做法假定同类对象共有这些属性，对类内变异、遮挡与不可见属性并不充分\scite{ket-lampert2009attributes}

属性预测要变成类别分数，还须说明概率模型。在理想化DAP模型中，类别$c$确定$\bm a_c$，即$p(\bm a\mid y=c)=\mathbf1\{\bm a=\bm a_c\}$，并假定给定属性后，类别与图像条件独立。记类别先验为$\pi_c=p(y=c)$，且$p(\bm a_c)>0$。先由贝叶斯公式求属性对应的类别概率，再对不可直接观察的属性求和，得到
\begin{equation}
  \begin{aligned}
    p(y=c\mid\bm a)
      &=\frac{\pi_c\,\mathbf 1\{\bm a=\bm a_c\}}{p(\bm a)},\\
    p(y=c\mid\bm x)
      &=\sum_{\bm a\in\{0,1\}^m}
          p(y=c\mid\bm a)\,p(\bm a\mid\bm x)\\
      &=\pi_c\,\frac{p(\bm a_c\mid\bm x)}{p(\bm a_c)}.
  \end{aligned}
  \label{eq:ket-transfer-attribute-link}
\end{equation}
类别描述固定以后，求和只剩$\bm a=\bm a_c$的一项。这一步尚未解决怎样从$m$个单属性预测器得到联合属性概率。为使比值可计算，再分别近似图像条件下的属性联合分布与属性先验，记$p_j=p(a_j=1)$：
\[
  \begin{aligned}
    p(\bm a_c\mid\bm x)
      &\approx\prod_{j=1}^{m}
        q_j(\bm x)^{a_{c,j}}[1-q_j(\bm x)]^{1-a_{c,j}},\\
    p(\bm a_c)
      &\approx\prod_{j=1}^{m}
        p_j^{a_{c,j}}(1-p_j)^{1-a_{c,j}}.
  \end{aligned}
\]
两个乘积分别假定条件属性独立与先验属性独立，前者不推出后者：混合不同图像后，属性仍可能相关。将它们代入式~\eqref{eq:ket-transfer-attribute-link}，便能只用$q_j$、$p_j$和目标描述计算类别分数：
\begin{equation}
  s_c(\bm x)=\pi_c
  \prod_{j=1}^{m}
  \frac{q_j(\bm x)^{a_{c,j}}
        [1-q_j(\bm x)]^{1-a_{c,j}}}
       {p_j^{a_{c,j}}(1-p_j)^{1-a_{c,j}}},
  \label{eq:ket-transfer-attribute-score}
\end{equation}
分子比较预测属性与类别描述，分母校正属性组合的常见程度。不过真实目标先验通常未知，DAP采用来源估计后，这个量应视为分类分数，不再是前述理想模型下的精确后验。Lampert等按来源类别属性表逐列取均值，不按各类图像数量加权。本节采用相同逐类口径，再加对称伪计数$\delta>0$避免分母为零：
\[
  \hat p_j=\frac{\sum_{c\in C_S}a_{c,j}}{|C_S|},
  \qquad
  p_j=\frac{\sum_{c\in C_S}a_{c,j}+\delta}{|C_S|+2\delta}.
\]
这里$C_S$非空，$\hat p_j$为未平滑的逐类均值，$p_j$为平滑后的实际计算值，必在$(0,1)$内；图像多的类别不会获得额外权重。对非空目标候选集，本节取均匀$\pi_c=1/|C_T|$，不以最终查询标签估计频率。连乘可改为累加对数，零预测概率则须按预定裁剪或极限约定处理。来源先验、属性独立性和预测器误差都可能偏离目标实际情况，所以将分数归一化也没有自动完成概率校准。

如果来源已经有类别概率模型，还可以先预测源类，再按源类属性表汇总属性。Lampert等将这条路径称为间接属性预测（indirect attribute prediction, IAP）：
\[
  p(a_j=1\mid\bm x)
  =\sum_{c\in C_S}a_{c,j}\,p(y=c\mid\bm x).
\]
这把$q_j$的学习换成对源类概率的复用，仍受源类与属性覆盖限制。两条路径都不能突破描述本身的区分能力：新种甲、乙的属性向量及先验相同，DAP分数就相同；两者仅在叶缘上不同，而照片没有叶片，观测也可能无法区分。增加描述词语只有在它们对应可观察、可学习的差异时才提供新依据。

预选属性词表限制了可表达的描述，连续图文表示则学习另一种比较接口。\term{对比语言图像预训练}（Contrastive Language--Image Pre-training, CLIP）训练图像与文本编码器，在对比目标中提高真实配对相对其他配对的相似度。来源文本由此提供视觉语义联系，而不只是固定类编号。对比目标怎样由配对数据构成，见第~\ref{sec:data-derived-knowledge-contrastive}节；迁移时使用已经训练好的两个编码器、共同空间及尺度约定，将目标描述编码为图像可比较的参照\scite{radford2021}

这个空间可编码多种表述，却不要求每个维度对应一个可解释属性，也不保证相似度能分辨否定、数量或细粒度形态。名称、属性和解释仍是目标知识，应记录提供者、模板选择所用验证材料及预训练覆盖。声称来源未见类，需要实际控制来源类别；网页预训练模型未在目标训练集上监督拟合，并不能证明它从未接触目标概念。

\subsection{参数初始化与任务先验的准备}
\label{sec:ket-transfer-initialization-prior}

少量新类图像也可以用来优化模型，但它们往往不足以约束所有参数。来源准备可以给更新一个起点，也可以规定任务参数的先验分布。初始化$\bm\theta$是一组具体数值；任务先验$p(\bm w_k\mid\bm\eta)$描述不同任务参数的可能取值，$\bm\eta$是共享超参数。先验均值可作初始化，却没有因此把方差、相关性或多峰结构交给更新过程。下面先讨论起点怎样按适应后损失学习，再讨论分布怎样由证据修正。

\subsubsection{参数初始化的学习}
\label{sec:ket-transfer-maml}

新任务若只允许几次梯度更新，直接预测源任务的最优参数未必是最合适的出发点。\term{模型无关元学习}（model-agnostic meta-learning, MAML）把支持更新后的查询损失作为准备目标，从历史任务中学习共享起点。它试图让少量支持反馈形成有用的任务模型，而非要求起点直接回答每个任务。“模型无关”限定于能够通过相应梯度过程训练的参数化模型，仍需选定结构、损失与任务协议\scite{finn2017}

对非空支持、查询集，记$L_{S_k}(\bm w)$和$L_{Q_k}(\bm w)$为平均损失。固定一批$M$个历史任务，先沿支持损失的负梯度做一步任务内更新，再平均各任务更新后的查询损失以调整共享起点：
\begin{align}
  \bm w_k^{(0)}&=\bm{\theta}, \nonumber\\
  \bm w_k^{(1)}&=\bm{\theta}
       -\alpha\nabla_{\bm{\theta}}L_{S_k}(\bm{\theta}),
       \label{eq:ket-transfer-maml-inner}\\
  J(\bm{\theta})&=\frac{1}{M}\sum_{k=1}^{M}
       L_{Q_k}(\bm w_k^{(1)}(\bm{\theta})), \nonumber\\
  \bm{\theta}&\leftarrow\bm{\theta}-\beta\nabla_{\bm{\theta}}J(\bm{\theta}).
       \label{eq:ket-transfer-maml-outer}
\end{align}
其中$\alpha>0$为任务内步长，$\beta>0$为跨任务步长，$\bm w_k^{(1)}$是支持更新后的任务参数。各任务从同一个当前$\bm\theta$独立复制，不能让下一任务继承上一任务的末态。外层以降低适应后的平均查询损失为目标，不保证每次更新都下降；保存的是经验证选出的共享起点，而非把各任务分类器依次拼成一个分类器。

查询损失对起点的依赖经过内层更新，不能直接用查询参数的梯度代替起点梯度。假设损失足够光滑，令$\bm H_k=\nabla^2L_{S_k}(\bm\theta)$，$\bm I$为与参数维数相符的单位矩阵。对内层式求导并应用链式法则，有
\begin{align}
  \frac{\partial\bm w_k^{(1)}}{\partial\bm{\theta}}
     &=\bm I-\alpha\bm H_k, \nonumber\\
  \nabla_{\bm{\theta}}J(\bm{\theta})
     &=\frac{1}{M}\sum_{k=1}^{M}
       (\bm I-\alpha\bm H_k)^{\top}
       \nabla_{\bm w}L_{Q_k}(\bm w_k^{(1)}).
       \label{eq:ket-transfer-maml-chain}
\end{align}
查询梯度在适应后参数处给出下降要求，雅可比$\bm I-\alpha\bm H_k$再把它换算成起点处的要求，这就是准备目标与直接查询训练的差异。二阶连续可微时Hessian对称；实现可用Hessian与向量的乘积，无须显式保存参数维数平方的矩阵。非光滑网络中的实际求导约定另由实现确定，不能据此宣称处处满足上述经典二阶条件。

多步内层更新会逐步改变起点到任务末态的映射。对预先确定、在此求导中不依赖参数的步长$\alpha_t$，令$\bm J_k^{(t)}=\partial\bm w_k^{(t)}/\partial\bm\theta$，则
\[
  \bm J_k^{(0)}=\bm I,\qquad
  \bm J_k^{(t+1)}
  =\bigl[\bm I-\alpha_t\nabla^2L_{S_k}(\bm w_k^{(t)})\bigr]
    \bm J_k^{(t)}.
\]
因此第$k$项的起点梯度为$(\bm J_k^{(T)})^\top\nabla L_{Q_k}(\bm w_k^{(T)})$，各步雅可比按实际更新顺序相乘；反向自动微分沿同一轨迹反传。普通混合任务训练直接降低共同参数上的合并数据损失，没有先形成任务解再评价的映射。两者即便使用相同骨干与优化器，也应在相同支持证据和更新预算后比较，更新前准确率回答不了起点是否利于适应。

算法~\ref{alg:ket-transfer-maml}将重置、适应和外层反馈串成完整训练过程。模型的任务输出接口须能表达历史与目标标签，初始参数的生成规则也须预先规定。目标使用只执行冻结方案中的支持更新，历史采样和跨任务求导均属于准备费用。

\begin{algorithm}[MAML的准备与目标适应]
  \label{alg:ket-transfer-maml}
  \begin{algorithmic}[1]
    \Require 历史训练任务、独立元验证任务；内层步数$T$与步长$\alpha_t$；外层步长$\beta$
    \Ensure 冻结的共享初始化及给定目标支持集的适应模型
    \State 初始化共享参数$\bm{\theta}$
    \While{训练预算尚未耗尽，且未满足元验证停止条件}
      \State 抽取$M$个历史任务，各自构造不重叠的$S_k,Q_k$
      \For{$k=1,\ldots,M$}
        \State 复制共享起点：$\bm w_k^{(0)}\gets\bm{\theta}$
        \For{$t=0,\ldots,T-1$}
          \State $\bm w_k^{(t+1)}\gets\bm w_k^{(t)}
            -\alpha_t\nabla L_{S_k}(\bm w_k^{(t)})$
        \EndFor
        \State 用$\bm w_k^{(T)}$计算查询损失$L_{Q_k}$
      \EndFor
      \State 沿内层更新轨迹求导，用平均查询梯度更新$\bm{\theta}$
    \EndWhile
    \State 按元验证结果选定并冻结初始化、步长与步数
    \State 新目标从该初始化出发，仅在$S_T$上执行约定的$T$步更新
    \State 对$Q_T$输入产生预测；隐藏答案只进入独立评分
  \end{algorithmic}
\end{algorithm}

\begin{example}[同一组分类任务中的起点与适应后损失]
  \label{ex:ket-transfer-maml-numeric}
  为了只比较起点与一步适应，构造一维花卉特征$z$，用$wz$作为二值标签$y\in\{-1,1\}$的预测分数，按$\operatorname{sign}(wz)$分类，并以平方损失训练。这不是MAML实测结果。任务1支持记录为$(1,1),(-1,-1)$，查询为$(2,1),(-2,-1)$；任务2支持为$(2,1),(-2,-1)$，查询为$(1,1),(-1,-1)$。记录代表独立样本，支持与查询幅度刻意不同，以展示准备目标对评价位置的依赖。

  对每组采用$\frac{1}{2n}\sum_i(wz_i-y_i)^2$，得到
  \[
    L_{S_1}(w)=\tfrac12(w-1)^2,\quad
    L_{S_2}(w)=\tfrac12(2w-1)^2,
  \]
  \[
    L_{Q_1}(w)=\tfrac12(2w-1)^2,\quad
    L_{Q_2}(w)=\tfrac12(w-1)^2.
  \]
  对两个支持损失求和并令导数为零，得普通混合训练起点$w_{\mathrm{pre}}=3/5$。固定同一步长$\alpha=1/4$，内层更新分别给出$w_1^{(1)}=3w/4+1/4$和$w_2^{(1)}=1/2$。将这两个末态代入查询损失并取平均，元目标为
  \[
    J(w)=\tfrac14(3w/2-1/2)^2+\tfrac1{16},
  \]
  其最小点是$w_{\mathrm{meta}}=1/3$。在普通起点$w=3/5$处，$\mathrm dJ/\mathrm dw=0.3$。任务2内层雅可比为$1-\alpha\cdot4=0$，说明一步后参数不再依赖起点，因此它的起点梯度为零。直接将适应后查询梯度作用于起点，会漏掉这个依赖关系。

  更新前，普通与元起点的平均查询损失依次为$0.05$与$5/36\approx0.1389$，普通起点较小；同样一步后，依次为$0.1025$与$0.0625$，元起点较小。分数符号在此均正确，准确率不能区分它们，平方分数损失却能。这个构造只说明固定更新预算下应评价适应后目标，不是实际花卉任务中MAML优于预训练的证据。
\end{example}

精确回传需要保存或重算内层计算图。\term{一阶MAML}（first-order MAML, FOMAML）将适应雅可比近似为单位矩阵，直接用适应后查询梯度调整共享起点，以省去经过支持梯度的二阶求导。历史任务采样、内层更新和查询反馈仍然存在；减少求导路径并没有取消元训练，也会改变梯度。

若改用任务内训练得到的位移，而非显式查询梯度，便得到Reptile的一阶过程。各历史任务独立从共享起点做若干步更新，获得$\bm w_k^{(T)}$，再以$\epsilon>0$控制向平均末态移动的幅度：
\begin{equation}
  \bm{\theta}\leftarrow\bm{\theta}
   +\epsilon\frac1M\sum_{k=1}^{M}
        \bigl(\bm w_k^{(T)}-\bm{\theta}\bigr).
  \label{eq:ket-transfer-reptile}
\end{equation}
Reptile不要求另取查询集计算外层梯度，方向来自任务内小批次更新的累计位移。一步普通SGD时，它退化为混合任务梯度方向；多步则在已更新参数上计算后续梯度，保留了任务内小批次的相互作用。FOMAML仍读显式查询梯度，故两者不是同一种去掉Hessian的算法。Reptile的独立元验证和最终目标查询隔离仍须保留\scite{ket-nichol2018reptile}

为看清位移中的相互作用，记同一任务两批损失在共同起点的梯度为$\bm g_1,\bm g_2$，第二批Hessian为$\bm H_2$。若第二批梯度可作带二阶余项的局部泰勒展开，固定小步长$\alpha$时有
\[
  \bm w^{(2)}-\bm{\theta}
    =-\alpha(\bm g_1+\bm g_2)
      +\alpha^2\bm H_2\bm g_1+O(\alpha^3).
\]
第二批梯度取在$\bm\theta-\alpha\bm g_1$，展开后产生$\bm H_2\bm g_1$项，而非仅在原起点相加两个梯度。Reptile实现仍只调用一阶梯度；这里用Hessian解释小步位移，没有要求算法额外执行二阶求导。

还可以控制内层范围，检查是否必须改变表示。将参数分为骨干$\bm\theta_b$与任务头$\bm\theta_h$，\term{几乎无内循环}（Almost No Inner Loop, ANIL）只在任务内更新头，外层仍训练骨干与头起点。查询损失直接依赖查询表示，也经支持表示依赖头适应，故精确ANIL没有在整个准备期间冻结骨干。原研究在MiniImageNet等指定网络与少样本任务中支持特征复用的解释；若新病害信息已被骨干丢失，目标端只调头仍无法恢复它\scite{ket-raghu2019anil}

\subsubsection{任务参数先验的学习}
\label{sec:ket-transfer-prior}

两张病害图像可能对应许多同样拟合的参数。一个起点没有描述这些候选的分布，而层次先验用共享超参数表达历史任务参数的联系，再让目标证据修正它。设$\bm\eta$为共享超参数，$\bm w_k$为任务参数；本节给定各支持输入，仅对标签建立似然。假定给定$\bm\eta$后任务参数独立，给定各$\bm w_k$后任务标签按似然独立，联合分解为
\begin{equation}
  p(\bm\eta,\{\bm w_k,S_k\}_{k=1}^{M})
   =p(\bm\eta)\prod_{k=1}^{M}
      p(\bm w_k\mid\bm\eta)
      \prod_{(\bm x_i,y_i)\in S_k}
      p(y_i\mid\bm x_i,\bm w_k).
  \label{eq:ket-transfer-hierarchical-model}
\end{equation}
式左的$S_k$是对支持标签及其已给定输入的简记，不是另建输入生成模型。要从历史数据选择先验参数，可积分掉未知$\bm w_k$，比较每项先验对该任务全部标签赋予的边缘概率。对预定的超参数候选范围，最大化这些概率的对数和：
\begin{equation}
  \hat{\bm\eta}
  =\argmax_{\bm\eta}\sum_{k=1}^{M}
     \log\int p(S_k\mid\bm w_k)\,
          p(\bm w_k\mid\bm\eta)\,\mathrm d\bm w_k.
  \label{eq:ket-transfer-empirical-bayes}
\end{equation}
积分保留每个先验允许的任务参数，而非先为各任务任选一个最优点；对数把独立任务概率的乘积转为求和。用历史数据估计先验超参数称为经验贝叶斯（empirical Bayes），另一个选择是保留$\bm\eta$的后验。两者都需要规定分布族、积分或优化近似与任务关系；一个任意正则系数不能据名称就成为已学到的先验。先验结构仍由独立历史验证任务选择，目标查询答案不参与上述估计。

准备固定$\hat{\bm\eta}$后，目标支持集的似然与先验相乘形成后验。若只需一个参数点，可在允许的任务参数范围内最大化后验，等价地最小化负对数似然与负对数先验之和：
\begin{align}
  p(\bm w_T\mid S_T,\hat{\bm\eta})
  &\propto p(S_T\mid\bm w_T)p(\bm w_T\mid\hat{\bm\eta}), \nonumber\\
  \hat{\bm w}_T
  &=\argmin_{\bm w}\left[
       -\sum_{(\bm x_i,y_i)\in S_T}\log p(y_i\mid\bm x_i,\bm w)
       -\log p(\bm w\mid\hat{\bm\eta})\right].
  \label{eq:ket-transfer-map}
\end{align}
第一行描述完整参数分布，第二行给出最大后验（maximum a posteriori, MAP）点；写$\argmin$不保证优化已找到全局最小值。要让参数不确定性进入预测，还需计算$\int p(y\mid\bm x,\bm w)p(\bm w\mid S_T,\hat{\bm\eta})\,\mathrm d\bm w$。它将各参数的输出按后验加权，不能以一个MAP点输出代替。

为计算先验强度的作用，暂将目标换成花卉连续测量的简化线性回归，不把类别编号当连续值。设$\bm w_k\in\mathbb R^d$服从$\mathcal N(\bm\mu,\sigma_w^2\bm I)$，$y_i=\bm z_i^\top\bm w_k+\varepsilon_i$，噪声独立服从$\mathcal N(0,\sigma_y^2)$，两方差严格为正。历史任务提供中心$\bm\mu$和差异尺度，噪声尺度亦按预定方案给定或估计。目标$n$个支持特征按行组成$\bm Z\in\mathbb R^{n\times d}$，观测组成$\bm y\in\mathbb R^n$。将似然与先验的二次项配方，后验协方差和均值为
\begin{align}
  \bm\Sigma_T
    &=\left(\sigma_y^{-2}\bm Z^\top\bm Z
                     +\sigma_w^{-2}\bm I\right)^{-1}, \nonumber\\
  \bm m_T
    &=\bm\Sigma_T\left(\sigma_y^{-2}\bm Z^\top\bm y
                     +\sigma_w^{-2}\bm\mu\right).
  \label{eq:ket-transfer-gaussian-posterior}
\end{align}
因$\sigma_w^{-2}\bm I$严格正定，矩阵可逆；$\bm\Sigma_T\in\mathbb R^{d\times d}$，$\bm m_T\in\mathbb R^d$。所得后验为$\mathcal N(\bm m_T,\bm\Sigma_T)$，此时均值恰好也是MAP点。支持数据不仅移动中心，还改变各参数方向的方差；单独保存起点$\bm\mu$无法表达这个更新后的分布。

\begin{example}[先验怎样改变少样本估计]
  \label{ex:ket-transfer-prior-numeric}
  构造一维$z_i=1$、噪声方差为1、先验$w_T\sim\mathcal N(0,1)$，支持观测为$2,4$。无先验的似然估计是均值3；按式~\eqref{eq:ket-transfer-gaussian-posterior}，后验精度即方差的倒数为$1+2=3$，均值为$6/3=2$，方差为$1/3$。仅将先验方差减为$1/4$，均值变为$6/(4+2)=1$，方差为$1/6$；将其增为4，均值为$6/(1/4+2)=8/3$，方差为$4/9$。强先验让本模型内的后验更集中，也把中心拉得更近；若真实目标参数偏离历史中心，这种集中不表示估计更准确。
\end{example}

Grant等给出MAML的层次贝叶斯解释，将有限步任务估计与跨任务准备联系到经验贝叶斯近似。联系依赖估计与近似方式，不是任意起点都已指定完整参数先验\scite{ket-grant2018hierarchical}

这个区别可在一维二次目标中核对。令负对数似然的有效项为$L(w)=h(w-\hat w)^2/2$，$h>0$是曲率，$\hat w$是其最小点。从$\mu$以固定步长$\alpha$做$T\ge1$步梯度下降，递推给出$w^{(T)}=(1-a)\hat w+a\mu$，其中$a=(1-\alpha h)^T$。若改为加入中心在$\mu$的二次先验惩罚，则
\[
  \argmin_w\left\{L(w)+\frac{\lambda}{2}(w-\mu)^2\right\}
    =\frac{h\hat w+\lambda\mu}{h+\lambda}.
\]
在$0<\alpha h<1$下，$0<a<1$，取$\lambda=ha/(1-a)>0$，这个MAP点恰等于早停参数，对应高斯先验$\mathcal N(\mu,\lambda^{-1})$。例如$h=1,\alpha=1/2,T=1$给出$\lambda=1$。等价对象仅是指定损失尺度下的参数点，先验精度还依赖曲率、步长与步数；非线性网络中的有限步梯度不因此成为全局MAP求解，更没有给出完整后验。多个当前任务怎样共同设计先验结构，见第~\ref{sec:ket-sharing-prior}节；这里讨论历史准备怎样支持未来估计。

\subsection{参数更新规则的准备}
\label{sec:ket-transfer-learned-optimizer}

有了起点，仍要决定每次支持梯度怎样改变参数。固定同一目标模型，步长、动量和缩放不同会得到不同轨迹；历史任务也可以用来训练这个更新规则。设任务参数为$\bm w_k^{(t)}$，更新器参数为$\bm\psi$，内部状态$\bm s_k^{(t)}$保存此前反馈。先按既定规则创建模型起点与必要的新头，再初始化更新器状态，随后用支持损失的梯度产生增量：
\begin{align}
  \bm s_k^{(0)}
    &=\operatorname{Init}_{\bm\psi}(\bm w_k^{(0)}), \nonumber\\
  \bm g_k^{(t)}
    &=\nabla_{\bm w}L_{S_k}^{(t)}(\bm w_k^{(t)}), \nonumber\\
  (\Delta\bm w_k^{(t)},\bm s_k^{(t+1)})
    &=U_{\bm\psi}(\bm g_k^{(t)},\bm w_k^{(t)},
             L_{S_k}^{(t)},\bm s_k^{(t)}), \nonumber\\
  \bm w_k^{(t+1)}
    &=\bm w_k^{(t)}+\Delta\bm w_k^{(t)}.
  \label{eq:ket-transfer-learned-update}
\end{align}
这里$L_{S_k}^{(t)}$为第$t$步所选支持小批次损失，$\bm g_k^{(t)}$为其参数梯度，$U_{\bm\psi}$输出与参数同形的$\Delta\bm w_k^{(t)}$及新状态。损失、梯度、参数和状态是否全部作为输入，由具体实现规定。目标反馈改变$\bm w_k$，历史训练得到的$\bm\psi$在目标端固定；前者承载当前预测，后者规定怎样学习，不能混为同一项迁移资产。

Andrychowicz等以循环网络实现学习更新器。其逐坐标设计让一个小型长短期记忆网络（long short-term memory, LSTM）共享权重，但每个参数坐标保留自己的历史状态，并输出该坐标增量。这样不必用一个全连接网络读取完整参数向量；梯度符号和幅度的预处理也帮助处理不同量级。历史梯度可以经状态影响本步增量，类似可学习的动量与步长，不过训练出这种行为不意味着继承手工优化器的凸问题收敛保证\scite{ket-andrychowicz2016-learning}

状态初始化也是保存接口的一部分。逐坐标$\operatorname{Init}_{\bm\psi}$可以置零或复制学得初态，每个新任务与独立重复都按同一规则重新建立。任务内支持小批次之间沿用状态，任务间不自动传递末态；如果要跨任务保留状态，须另行声明连续适应协议。更新器只能读取许可的支持反馈，最终查询答案不进入状态或参数更新。

要训练$\bm\psi$，应评价它生成的轨迹，而不是仅拟合某个手工步长。将更新展开$T$步，以非负权重$\omega_t$规定哪些时刻需要好的查询预测，对$M$个历史任务取平均：
\begin{equation}
  J(\bm\psi)=\frac1M\sum_{k=1}^{M}
       \sum_{t=1}^{T}\omega_t L_{Q_k}(\bm w_k^{(t)}(\bm\psi)),
       \qquad \omega_t\geq0.
  \label{eq:ket-transfer-optimizer-objective}
\end{equation}
只令$\omega_T$非零，评价的是终点；多个非零权重还约束中途表现，权重须预先规定并至少有一个为正。少样本学习用独立查询评价适应后的泛化，纯函数优化也可直接用所优化的函数值，不能将两者混成同一目标。反向传播沿任务参数与更新器状态训练$\bm\psi$；截断轨迹或停止对梯度输入再求导，会省计算也会改变得到的梯度，不再是原展开目标的精确导数。

Ravi与Larochelle把少样本分类器参数与门控循环状态相联系，学习保留与更新的幅度。核心写成$\bm w^{(t+1)}=\bm f^{(t)}\odot\bm w^{(t)}-\bm i^{(t)}\odot\bm g^{(t)}$，其中$\odot$为逐坐标乘法，$\bm f^{(t)}$为保留门，$\bm i^{(t)}$为更新门，维度均与$\bm w$相同。损失、梯度和循环历史影响这些门，初态还可参与训练，从而联合准备起点与规则。查询标签用于历史外层评价，不交给本任务更新器\scite{ket-ravi2017optimization}

VeLO进一步用更广的神经网络优化任务训练更新器，并利用任务与训练进度信息调整优化行为。论文中无需逐任务搜索学习率的使用方式，仍依赖更新器准备、模型选择和执行成本，不能解释为全过程免调参。换成更大模型、更长轨迹、不同损失或梯度尺度时，要另测外推；短轨迹终点较低，不保证长时间训练稳定\scite{ket-metz2022velo}

费用也分为共享与任务状态。保存更新器需$|\bm\psi|$个参数；若任务有$P$个参数，每坐标循环状态宽度为$h$，状态另占$O(Ph)$存储，每步还增加更新器前向计算。少做若干次学习率搜索可节省目标端费用，却未必抵消历史训练、共享存储和逐步执行开销，应按实际任务数量核算。

\subsection{任务条件化机制的准备}
\label{sec:ket-transfer-conditioning-preparation}

几张支持图像也可以直接进入预测计算，而不先用于梯度更新。系统读取输入答案对，形成当前任务状态，再用这个状态回答查询。\term{任务条件化}（task conditioning）指预测显式依赖本次任务证据：共享参数规定怎样读，任务状态随证据改变。与学习更新器不同，这个过程可以不输出一套优化后的任务权重；准备阶段需要训练的是证据到条件预测的映射。

用固定语义的回归或分类任务说明这种映射。\term{条件神经过程}（Conditional Neural Process, CNP）将每个支持输入答案对编码成摘要，再由摘要与查询输入共同产生分布。对非空$S_k$，编码器$e_{\bm\phi}$输出$d_r$维向量，简单均值聚合为
\begin{align}
  \bm r_k&=\frac1{|S_k|}\sum_{(\bm x_i,y_i)\in S_k}
                   e_{\bm\phi}(\bm x_i,y_i), \nonumber\\
  p_{\bm\phi}(y\mid\bm x,S_k)
      &=p_{\bm\phi}(y\mid\bm x,\bm r_k).
  \label{eq:ket-transfer-cnp}
\end{align}
平均使$\bm r_k\in\mathbb R^{d_r}$不依赖支持排列，代价是可能丢掉计数与多峰细节；需要这些信息时须扩展聚合接口。回归解码器可输出均值和正方差，分类解码器可输出预定标签空间的概率。若未来类别数或标签含义可变，不能只沿用一个固定分类输出，必须另外规定标签编码与输出接口。共享$\bm\phi$由历史条件负对数似然训练，目标端冻结它，只从新支持计算$\bm r_T$。

原始CNP训练可从一组观测抽上下文，再预测包含上下文和未观察点的目标集。本章若评价支持之外的泛化，应另报告隔离$Q_k$上的损失，上下文重构不能替代它。编码器和解码器规模固定时，简单均值方案处理$n$个支持与$q$个查询约需$O(n+q)$次点级计算，固定支持摘要可缓存；这项复杂度不自动适用于全点注意力等扩展。

基本因子化解码器在固定支持条件下独立产生查询输出，因而没有表示完整的查询相关性。不同支持集得到的条件分布，也不保证均来自同一随机函数先验的贝叶斯条件化。输出均值和方差仅规定了模型的预测形式，既不自动给出共同先验解释，也不证明方差已校准，后者仍需真实查询检验\scite{ket-garnelo2018cnp}

若任务输入输出本来采用序列接口，可以将支持记录直接放入序列。MetaICL从历史任务抽取支持输入答案对与独立查询，将支持和查询输入拼接，以查询答案的负对数概率训练读取示范的能力。用可选说明扩展这一接口，准备目标写为
\begin{equation}
  \min_{\bm\phi}
  \mathbb E_k\left[
    -\log p_{\bm\phi}
       (y\mid D_k,(\bm x_1,y_1),\ldots,(\bm x_n,y_n),\bm x)
  \right],
  \label{eq:ket-transfer-metaicl}
\end{equation}
式中期望包括历史任务及其支持、查询采样，最小化在预定模型参数范围内进行。末尾$(\bm x,y)$为查询，$y$只进入损失，不作为预测自身的输入；序列答案的对数概率由词元条件概率求和。$D_k$是本书为“说明加示范”保留的可选扩展，原始MetaICL核心是示范与查询，并不要求说明。准备后，未见任务通过读取新示例改变预测而不更新权重，这是显式训练上下文适应的路径\scite{ket-min2022metaicl}

也可以直接取已有语言模型作为条件预测器。GPT-3的少样本使用把示例放入提示，而不逐目标梯度微调；其主要准备是大规模语言建模，不是上述显式任务训练。因此输入有示例，只证明当前接口使用了示例，不能据此判断训练采用元学习，更不能把收益全归给正确示范标签。应记录来源模型、任务覆盖和目标证据格式，并用控制实验检查证据的实际作用\scite{brown2020}

无论准备的是表示、语义编码器、先验、更新器还是条件模型，都须保存使用所需版本与初始化规则。它们的训练与验证进入式~\eqref{eq:ket-reuse-amortization}的准备费用$C_{\mathrm{prep}}$。要判断多次使用是否抵消这项投入，可将来源选择并入本章$C_{\mathrm{adapt},k}$，以$C_{\mathrm{predict},k}$记单次预测费用，对$K$个实际目标任务累加：
\begin{equation}
  C_{\mathrm{total}}
    =C_{\mathrm{prep}}
     +\sum_{k=1}^{K}
       \bigl(C_{\mathrm{adapt},k}+q_kC_{\mathrm{predict},k}\bigr),
  \label{eq:ket-transfer-total-cost}
\end{equation}
这里$K$仅表示本段实际任务数，$q_k$是任务$k$的请求次数；若逐次费用不同，应按真实请求逐项累加，而非强用一个常数。检索与状态管理依发生阶段归入适应或预测，所有计算费用须用相同单位，人工和存储另列或按明确单价折算。第~\ref{sec:ket-reuse-resource-benefits}节要求计入失败尝试、匹配无迁移对照，并区分生命周期总投入与利用已有模型的增量决策。较低历史准备损失还不是未来任务收益的证据。

\section{目标任务的适应}
\label{sec:ket-transfer-adaptation}

现在固定一个已经交付的来源模型，目标端怎样使用新种甲、乙的证据？描述或支持表示可以形成显式参照；说明和示范可以成为条件模型的上下文；监督也可以通过优化写入任务参数。下面按直接改变的对象比较这三种操作。它们不是必经阶段，可以组合，例如先建立文本参照，再用少量标注拟合校正头；不更新网络权重的前两种方式也仍然使用目标知识。

\subsection{预测参照的构建}
\label{sec:ket-transfer-references}

只给新种甲、乙的描述时，先限定预测只能在这两个指定未见类中选择。这是传统零样本学习（zero-shot learning, ZSL）的候选条件；已见类也可能出现在测试中时，则进入广义零样本学习（generalized zero-shot learning, GZSL）。若本段$C_T$专指未见类，候选集扩为$C_S\cup C_T$，须分别检查两部分错误。来源已见类的训练优势、候选数量与归一化都改变了，原ZSL阈值不能原样继承。

描述足以区分类别时，DAP直接使用式~\eqref{eq:ket-transfer-attribute-score}；CLIP类模型则将文本参照和查询图像编码为非零向量，再分别归一化成单位向量$\bm t_c,\bm z\in\mathbb R^d$。为在当前$N$个候选类中比较，采用点积相似度并以温度$\tau>0$归一化：
\begin{equation}
  p(y=c\mid\bm x,D_T)
  =\frac{\exp(\bm z^\top\bm t_c/\tau)}
         {\sum_{c'=1}^{N}\exp(\bm z^\top\bm t_{c'}/\tau)}.
  \label{eq:ket-transfer-text-reference}
\end{equation}
单位向量的点积就是余弦相似度，值越大表示在已学空间中越近。参照内容来自描述，比较机制来自来源图文训练；得到的是候选集内的模型概率，不自动是校准后验。加入新候选会改变分母与原类别概率，即使编码器未改。使用多个模板时，可按预定规则平均编码并重新归一化，但模板选择的验证材料和编码费用仍须记账。

\begin{example}[描述、支持样本与原型的比较计算]
  \label{ex:ket-transfer-reference-numeric}
  构造两类描述与向量，分别计算语义和图像支持的作用，数值不代表实际植物属性或模型实测结果。教学属性依次为“花瓣细长”“叶缘有锯齿”，甲类描述$(1,0)$，乙类$(0,1)$。查询属性预测$(0.8,0.3)$，两类先验均为$1/2$；来源属性表各列有一半类别取1，逐类均值与对称平滑均给出$1/2$。DAP公共因子消去后，甲、乙分数正比于$0.56$与$0.06$，分别归一化为$28/31\approx0.9032$、$3/31\approx0.0968$。

  另一组共同空间中，单位文本参照为$\bm t_{\text{甲}}=(1,0)$、$\bm t_{\text{乙}}=(0,1)$，单位图像表示为$\bm z=(0.8,0.6)$。取$\tau=0.2$，两类相似度经缩放为4和3，甲类概率为$e^4/(e^4+e^3)\approx0.7311$。两次都选甲，概率却来自不同假设与表示，不能比较数值大小来判断哪种方法更可靠。

  再让目标提供每类两张标注图像，以另一个冻结视觉编码器得到甲类$(0,0),(0,2)$、乙类$(2,0),(2,2)$，查询为$(0.5,1)$。原型为$(0,1)$、$(2,1)$，查询平方距离为$0.25,2.25$，温度1给甲概率$1/(1+e^{-2})\approx0.8808$。同表示上逐例匹配的距离为$1.25,1.25,3.25,3.25$，相同负平方分数求和也得到这个概率。这次相等来自对称点的共同偏移，不是多样本时两规则恒等，更不与前两种证据条件等价。
\end{example}

要检验均值丢掉了什么，换成一维教学支持：甲为$-2,2$，乙为$0.5,0.7$，查询为2。原型0与0.6的平方距离是4与1.96，故单原型选乙；温度1的逐例匹配给甲总权重$1+e^{-16}$，乙总权重$e^{-2.25}+e^{-1.69}$，故选甲。甲的两端模式平均到了没有支持点的位置，说明减少参照改变了判别规则。保留逐例或增加类内原型可表达更多模式，但极少支持也可能不足以估计这些模式。

目标存在私有类别时，候选内比较还缺退出机制。部分重叠任务须识别共享$C_{\cap}$，开放集任务则须允许未知输入不接受已知参照。令非空$C_{\mathrm{known}}$为当前已知候选，以最大类别概率作为一个简单接受分数：越大越倾向于接受已知类，达到预定阈值$\gamma$才分类，
\begin{equation}
  \hat y(\bm x)=
  \begin{cases}
    \argmax_{c\in C_{\mathrm{known}}}p(c\mid\bm x),
       & \max_{c\in C_{\mathrm{known}}}p(c\mid\bm x)\geq\gamma,\\
    \mathrm{unknown}, & \text{否则}.
  \end{cases}
  \label{eq:ket-transfer-rejection}
\end{equation}
这里$\gamma\in[0,1]$，并列最大类按固定规则处理。阈值依独立验证、许可的已知误拒率或成本约束确定，不可用最终未知测试答案选到最佳。softmax始终归一，即使全部已知类都不适合，也可能给一个类很高的概率，所以这只是待验证的拒识规则。距离阈值、未知验证类或专门拒识模型提供另一类依据，也增加信息与假设；仅有已知类验证，最多约束相应误拒，不能保证任意未知类检出。

分类参照只回答候选类别问题，不能继承为任意输出任务。检测可在候选区域上比较参照，但还须产生区域、恢复位置并去除重复；分割须保留空间分辨率，生成须接入序列解码条件。另一方面，冻结特征上拟合线性头虽然也由支持建立分类规则，却用优化形成任务参数，属于第~\ref{sec:ket-transfer-parameter-adaptation}节。直接均值与参数拟合的区别是如何形成预测器，不是有没有使用标签。

\subsection{任务上下文的组织}
\label{sec:ket-transfer-context}

同一张照片既可用于物种识别，也可用于病害复查；条件模型必须从当前证据知道要回答哪个问题。\term{上下文学习}（in-context learning, ICL）指模型利用当前输入中的示例形成本次任务行为。本节讨论冻结共享权重的使用方式，变化发生在激活、注意力读取或任务状态中，不自动成为下次会话仍保存的参数知识。第~\ref{sec:ket-transfer-conditioning-preparation}节说明怎样准备这种接口，这里检查目标怎样组织可读证据。

例如指令规定“按观察记录输出病害复查等级”，教学示范给出“无斑点且边缘完整$\rightarrow$低”“有扩散斑点$\rightarrow$高”，最后放入无答案的新记录。说明定义责任，示范展示对应，查询是待预测对象。把指令改成物种识别，原示范就没有提供新目标的映射。这里使用的是文本观察记录；要检查真实图片还需图像编码接口，文本模型不能直接读取未编码图片。序列边界必须明确，查询答案只保留在评价端。

若多个规则都可能符合当前证据，可以显式计算它们的权重。\term{任务后验推断}（task posterior inference）为潜在规则或任务变量$u$规定先验及证据似然，再对支持与说明得到的后验加权预测。假设$u$给出标签机制，且查询输入不额外改变规则后验，则
\begin{equation}
  p(y\mid\bm x,S_T,D_T)
  =\int p(y\mid\bm x,u)\,
          p(u\mid S_T,D_T)\,\mathrm du.
  \label{eq:ket-transfer-task-posterior}
\end{equation}
积分覆盖所设规则或参数范围，不是只选一条最高权重规则；离散$u$改用求和。例如候选中可包含“扩散斑点出现就判高等级”。若任务身份还决定输入分布，当前查询$\bm x$也可提供任务信息，后验须改为$p(u\mid\bm x,S_T,D_T)$。允许联合使用其他查询时，还应按直推协议把它们纳入条件，不能不改条件就扩大信息范围。

不显式建后验时，模型也可以从外部支持记录选择示范或计算读取向量，这称为记忆读取；直接把描述与示范输入网络，则是上下文条件化。任务后验强调概率含义，记忆读取强调证据访问，条件化强调预测依赖，三者可同时存在。例如网络状态可以编码近似后验，注意力也可读取外部记录。核算成本与撤销证据时，仍要列清实际保存的是后验参数、样本记录还是激活，以及缓存多久。

Xie等在文档由混合隐马尔可夫模型生成的设定中，分析从语言建模到文档级潜在概念推断的联系。隐马尔可夫模型让不可直接观察的状态影响序列生成；相应论证还限制模型拟合、提示与预训练分布差异及概念可区分性。上述条件使后验解释可分析，却没有证明任意实际语言模型都学到正确先验或执行精确积分\scite{xie2022bayesianicl}

另一种可核对的解释是，前向读取可以得到与一步参数更新相同的预测。为限定对象，取$d$维特征的线性回归$\hat y=\bm w^\top\bm z$，对$n>0$个支持使用平方损失$L(\bm w)=\frac1{2n}\sum_i(\bm w^\top\bm z_i-y_i)^2$。求其梯度、做步长$\alpha$的一步更新，再与查询$\bm z$相乘，得到
\begin{equation}
  \hat y(\bm z)
  =\bm w^{(0)\top}\bm z
   -\frac{\alpha}{n}\sum_{i=1}^{n}
     (\bm w^{(0)\top}\bm z_i-y_i)\,\bm z_i^\top\bm z.
  \label{eq:ket-transfer-context-gd}
\end{equation}
支持残差$r_i=\bm w^{(0)\top}\bm z_i-y_i$决定每条记录应带来的修正。取查询向量$\bm q=\bm z$、键$\bm k_i=\bm z_i$、值$v_i=-\alpha r_i/n$，未归一化读取$\sum_i(\bm q^\top\bm k_i)v_i$恰好等于式中的修正项。这里不做softmax，不除权重和，因此可以有负权重，也不同于概率注意力。前向计算于是可产生一步梯度更新后的同一预测，而不持久写入$\bm w^{(1)}$。von Oswald等给出特定编码下的线性注意力构造并研究简单回归训练；这种代数对应不能直接推广成任意Transformer执行梯度下降的结论\scite{ket-vonoswald2023gradient}

这些构造解释某些机制，却不能判断当前模型到底利用了什么。Min等在所研究分类和多项选择任务中，用标签替换等对照观察到标签空间、输入分布与格式的作用，不能据此断言任意新映射都不需要正确标签\scite{ket-min2022demonstrations}表~\ref{tab:ket-transfer-context-controls}给出对应检查：固定查询和基本预算，每次只改变一个因素，观察同一目标指标怎样变化。

\begin{table*}[tbp]
  \centering
  \small
  \begin{tabularx}{\linewidth}{p{30mm}XX}
    \toprule
    受控变化 & 保持不变的对象 & 希望区分的解释 \\
    \midrule
    将示范标签随机替换 & 支持输入、标签词表、序列长度 & 是否依赖正确输入答案对应，还是主要利用格式与标签范围 \\
    将类别名改为新符号并给出一致映射 & 输入和正确对应关系 & 是否学会当前映射，还是沿用预训练中的类别名称习惯 \\
    交换示范顺序 & 全部输入答案对与查询 & 顺序或邻近位置是否主导结果 \\
    仅改变格式和分隔符 & 任务语义与所有示范内容 & 是否对表示接口敏感，而不是知识不足 \\
    更换任务指令或移除示范 & 查询输入与允许信息来源 & 任务说明和示范各自贡献什么，二者冲突时依赖谁 \\
    \bottomrule
  \end{tabularx}
  \caption{检查上下文依赖的单因素对照。每行独立执行，固定查询并记录改变因素之外的证据与长度；改名须在支持和正确评价映射中一致，错误标签实验则有意破坏对应。对照本身不预设结果，提示选择只使用验证任务，最终查询答案不参与选择。}
  \label{tab:ket-transfer-context-controls}
\end{table*}

目标端检索可以按查询相似性选择示范，但库中只能有许可记录，不能取回查询自身或按评价单位应隔离的同源答案。较长序列增加编码与生成费用；因果语言模型能缓存固定前缀的键值，缓存却占存储，示范变化还须重算相关部分。跨请求记忆要明确任务切换与清空规则，避免测试任务相互污染。这里仅处理被动证据；强化学习的任务推断还需主动选择动作获得反馈，并把探索费用及回报纳入评价，相应迁移条件可参考\BookChapterRef{chap:rl-transfer}{强化学习的迁移}。

\subsection{模型参数的调整}
\label{sec:ket-transfer-parameter-adaptation}

支持标签还可以通过优化形成持久任务参数。为了不把“训练多少”与“参数从哪来”混淆，先分开两项选择：允许多少来源已有参数更新，以及是否另建任务参数。前者分为不更新、部分更新、全部更新，后者分为不新增、新增。只训练新头是不更新已有参数但新增，全量微调加新头是全部更新并新增，LoRA加解冻归一化量则是部分更新并新增。表中的软提示是由目标损失训练的一组连续输入向量，而不是现成的可读文本；因此即使来源参数不更新，它仍然属于新增任务参数。下面按全量、已有部分和附加参数展开机制，这三个重点不构成互斥类别。

\begin{table}[htbp]
  \centering
  \small
  \begin{tabularx}{\linewidth}{p{28mm}XX}
    \toprule
    已有参数更新范围 & 不新增任务参数 & 新增任务参数 \\
    \midrule
    不更新 & 直接使用或上下文条件化已有模型 & 新建头、适配器、LoRA、软提示 \\
    部分更新 & 调整已有头、偏置、归一化量或部分层 & 新建头并解冻高层；LoRA并调整已有归一化量 \\
    全部更新 & 输出语义兼容时全量微调原模型 & 新任务头与全部来源参数共同微调 \\
    \bottomrule
  \end{tabularx}
  \caption{分别记录已有参数的更新范围与任务参数的新增。每格是一种参数来源组合，表中方法名仅代表常用配置，具体冻结与更新仍须按实现说明。新增参数少不保证来源全部冻结，也不保证反向传播路径短。}
  \label{tab:ket-transfer-parameter-axes}
\end{table}

参数范围也不决定来源怎样准备。普通预训练、MAML、先验估计和学习更新器，可以在接口相容时与不同范围结合。使用学习更新器，须按式~\eqref{eq:ket-transfer-learned-update}重建当前任务状态，并让规则只作用于选定参数。比较从相同支持证据出发，记录初始化、更新集合、步长与停止条件；配置依许可的独立验证确定，不能用最终标签选择。

\subsubsection{全量参数的调整}
\label{sec:ket-transfer-full-finetuning}

新物种若需要改变来源表示，\term{全量微调}（full fine-tuning）允许全部来源参数随目标监督调整。输出语义兼容时，目标从$\bm\theta_T^{(0)}=\bm\theta_S$开始；需要新类别或新任务头时，先规定头的结构与初始化，再与来源骨干一起训练。此时来源参数全量调整，新头另行学习，不能把未存在于来源的头也称为继承参数。

固定单标签新物种分类，用$p_{\bm\theta_T}(y\mid\bm x)$表示目标类别概率。要提高非空支持集正确类别的概率，同时限制允许的参数变化，可在预定任务模型范围内最小化
\begin{equation}
  L_T(\bm{\theta}_T)
   =-\frac1{|S_T|}\sum_{(\bm x_i,y_i)\in S_T}
       \log p_{\bm{\theta}_T}(y_i\mid\bm x_i)
       +\lambda\Omega(\bm{\theta}_T),
  \label{eq:ket-transfer-finetuning-loss}
\end{equation}
第一项是平均交叉熵，$\Omega$是预先规定的正则项，$\lambda\ge0$控制其权重；若取自先验，应交代与负对数先验及损失尺度的对应。按训练小批次执行$T$步优化，依据许可验证保存检查点。没有目标验证标签，就使用历史验证任务预先确定的配置，不能临时以查询集选步数。改为检测还须更换位置与匹配损失，单标签交叉熵本身不评价对象坐标。

要判断来源起点有无帮助，应固定目标结构、支持、增强与测试，比较来源初始化和随机初始化在相同预算下的适应曲线。充分训练后的最终风险也可比较，但它回答的不是少量更新预算问题。全量调整允许改变表示，也可能拟合少样本中偶然的背景关系；是否出现这种过拟合，取决于独立查询表现，不能仅由参数数量判定。

来源有$P$个参数时，全量训练通常为它们计算梯度，Adam等优化器还维护参数对应的状态，并处理中间激活。逐任务保存完整副本，权重存储约为$P$个参数，另加新头及必要状态。步数少不代表显存或总费用低于参照构建。目标风险下降还只说明当前能力，原物种是否仍可识别须另测。

新类别同时伴随拍摄或观测变化时，可结合第~\ref{sec:ket-update-weighting}节的校正及式~\eqref{eq:ket-update-crossmodal}的配对监督，但须重新核对标签联系、目标覆盖和访问权限。所需无标签批次、配对观测或旧数据并非每个协议都提供。保留旧能力的责任还需加入第~\ref{sec:ket-update-retention-validation}节的约束与评价，一次目标适应不能代替长期维护检查。

\subsubsection{已有部分参数的调整}
\label{sec:ket-transfer-existing-parameters}

目标仍识别来源的月季与鸢尾，但背景改变时，可以先尝试只调整原有头或部分层，而非开放全部参数。将来源参数分为冻结$\bm\theta_F$与可训练$\bm\theta_A$，固定来源冻结值，仅对允许子集最小化已定义的目标损失：
\[
  \min_{\bm{\theta}_A}
  L_T(\bm{\theta}_F=\bm{\theta}_{S,F},\bm{\theta}_A).
\]
候选范围由选定的层或坐标确定，更新只作用于$\bm\theta_A$，其余参数保持来源值。原有头、最高层、已有偏置和归一化量都是可选子集。“已有”按来源中是否存在判断，而非按训练占比或部署是否另存副本判断。

在这个语义兼容任务中，只调头改变固定表示中的边界，开放最高层还能改变高层特征组合，全量调整则允许早期特征变化。若后来改成依赖叶片细斑的病害任务，而最终表示已丢掉该差异，只读最终表示的头就无法补回。开放范围应结合目标所需信息与支持证据约束能力选择，不能用“层越多越好”代替同任务验证。

冻结仅规定权重不变，并不删除计算。若只训练顶部头，骨干与预处理均固定，可缓存支持特征供重复拟合，但新查询仍需骨干前向。可训练部件在中间时，损失梯度还须经过后续冻结层回传；不求这些层的权重梯度，仍可能需求输入梯度与激活。因此可训练参数比例下降，不自动给出同比例显存或推理节省。

ANIL在目标端只适应任务头，却使用跨任务准备过的骨干，说明部分更新与元学习可以组合。如果目标类别数或病害语义要求重新创建$\bm W_T,\bm b_T$，这些新参数就不再属于已有头更新。应列出复制、替换、初始化与冻结的操作；新头同时解冻骨干末层，是附加参数与已有部分更新的组合。

\subsubsection{附加任务参数的学习}
\label{sec:ket-transfer-added-parameters}

原类别集合不含新物种时，最直接的新增是任务头。固定来源$d$维编码器$h_{\bm\theta_S}$，创建$N$类权重$\bm W_T\in\mathbb R^{N\times d}$和偏置$\bm b_T\in\mathbb R^N$，将线性分数归一化为
\begin{equation}
  \bm p(\bm x)
   =\operatorname{softmax}(\bm W_T h_{\bm{\theta}_S}(\bm x)+\bm b_T).
  \label{eq:ket-transfer-new-head}
\end{equation}
第$c$项对应当前类别$c$的概率，目标训练的新增量共$N(d+1)$个。记$\bm z=h_{\bm\theta_S}(\bm x)$，真实标签独热向量为$\bm e_y$；单样本交叉熵对权重与偏置的梯度是$(\bm p-\bm e_y)\bm z^\top$与$\bm p-\bm e_y$，因而支持标签直接改变新头。平方距离原型也可改写成受约束线性分数，却由支持均值形成，不是这样独立优化权重。

只加头不能恢复固定最终表示丢掉的信息。若要在底座权重不变时调整中间计算，可以插入任务模块。\term{适配器}（adapter）让隐藏表示经过一个小型变换，并把增量加回原表示。典型瓶颈残差对$\bm h\in\mathbb R^d$计算
\begin{equation}
  \bm h'=\bm h+
     \bm W_{\mathrm{up}}\,
       \sigma(\bm W_{\mathrm{down}}\bm h+\bm b_{\mathrm{down}})
       +\bm b_{\mathrm{up}},
  \label{eq:ket-transfer-adapter}
\end{equation}
降维矩阵$\bm W_{\mathrm{down}}\in\mathbb R^{r\times d}$、升维矩阵$\bm W_{\mathrm{up}}\in\mathbb R^{d\times r}$，偏置分别为$r$维和$d$维，$\sigma$是规定的非线性激活，通常$r\ll d$。每个模块新增$2dr+r+d$个参数，小增量初始化让起始计算接近底座。目标损失训练新增量，部署保留模块及底座；每输入位置两次乘法约需$2dr$次乘加，另加激活、偏置与残差，多层和长序列会累积开销。含非线性的模块通常不能吸收为原线性矩阵的固定增量。

Houlsby等的配置在每层注意力输出投影后与前馈子层后各插入串行适配器，处于对应底座残差相加和归一化之前；适配器内部残差与底座残差应分别识别。原配置还训练归一化和任务输出部分，所以不能只数瓶颈矩阵。复制任务专有归一化量时须计入新增状态，直接调整来源已有归一化量则是部分已有参数更新\scite{ket-houlsby2019-adapters}

若希望调整某个线性映射而避免保存同大小的完整增量，可以将增量写成两个瘦矩阵之积。\term{低秩适应}（Low-Rank Adaptation, LoRA）固定$\bm W_0\in\mathbb R^{d_{\mathrm{out}}\times d_{\mathrm{in}}}$，新增$\bm A\in\mathbb R^{r\times d_{\mathrm{in}}}$、$\bm B\in\mathbb R^{d_{\mathrm{out}}\times r}$。对$\bm h\in\mathbb R^{d_{\mathrm{in}}}$，两条线性支路相加为
\begin{equation}
  \bm h'=\bm W_0\bm h+s\bm B\bm A\bm h,
  \label{eq:ket-transfer-lora}
\end{equation}
其中$s$是固定缩放，原形式常取$s=\alpha_{\mathrm L}/r$，$\alpha_{\mathrm L}$与本章梯度步长不同。新增$r(d_{\mathrm{in}}+d_{\mathrm{out}})$个参数，增量秩不超过$r$；这约束的是本次权重变化，不是证明知识具有低秩语义，也不使底座$\bm W_0$成为低秩。与非线性适配器相比，LoRA的并行增量保持为线性映射\scite{ket-hu2021lora}

训练仍需能从初态产生梯度。LoRA常随机初始化$\bm A$并将$\bm B$置零，使增量起始为零。记$\bm\delta=\partial L/\partial\bm h'\in\mathbb R^{d_{\mathrm{out}}}$，按乘法链式法则，单输入梯度为
\[
  \frac{\partial L}{\partial\bm B}
       =s\bm\delta(\bm A\bm h)^\top,\qquad
  \frac{\partial L}{\partial\bm A}
       =s\bm B^\top\bm\delta\,\bm h^\top.
\]
首步$\bm B=0$使$\bm A$梯度为零，但只要相应输入与损失梯度非退化，$\bm B$可以开始更新。若两因子都置零，两个梯度均为零，普通梯度过程无法启动这项增量。初始函数不变与参数可学习，必须同时检查。

线性增量允许在任务固定、格式兼容时预先形成$\bm W=\bm W_0+s\bm B\bm A$，之后只运行同形状矩阵。合并约需$O(d_{\mathrm{out}}d_{\mathrm{in}}r)$运算，并非免费。动态任务服务也可保留共享底座及各任务因子，未合并分支每输入增加$O(r(d_{\mathrm{in}}+d_{\mathrm{out}}))$计算。量化、并发与切换频率影响部署选择；“可合并”只消除指定形态下的额外预测分支，不消除加载和切换费用。

附加量还可进入输入端，而非隐藏层或权重矩阵。\term{软提示}（soft prompt）是一组由目标损失学得的连续输入向量。原输入嵌入$\bm E\in\mathbb R^{\ell\times d}$有$\ell$个位置，新增$\bm P_T\in\mathbb R^{m\times d}$，按行拼接$[\bm P_T;\bm E]$交给冻结模型，以目标答案负对数似然训练$md$个提示参数。它们通常不是可直接阅读的词语。Lester等的提示调优（prompt tuning）主要在输入嵌入处新增向量，不应与逐层键值前缀采用同一参数计数\scite{ket-lester2021prompt}

软提示也要用监督与反向传播形成。底座冻结时，求输入提示梯度仍需穿过模型层；部署保存提示与底座，长度从$\ell$增为$\ell+m$。稠密自注意力的分数计算从$O(\ell^2d)$增为$O((\ell+m)^2d)$，多出的主要项为$O((2\ell m+m^2)d)$，还会占用上下文与缓存。参数少不免逐次费用。文本示范直接作为上下文输入，没有这一步连续参数训练，证据形式与持久状态也不同。

图~\ref{fig:ket-transfer-added-locations}将这四种机制放在前向路径上比较。应观察新增量读入哪种对象：头读取最终表示，适配器变换中间表示，LoRA并行修改线性映射，软提示改变输入条件。这个位置说明可调整的计算不同，不能仅据位置给出同任务风险排序。

\begin{figure*}[htbp]
  \centering
  % Four forward-computation mechanisms; all dimensions are in millimetres.
% No external libraries beyond arrows.meta, positioning and calc are needed.
\begin{tikzpicture}[
  x=1mm,y=1mm,
  font=\figurefont\footnotesize,
  text=BookInk,
  >={Latex[length=2mm,width=1.4mm]},
  frozen/.style={draw=BookMuted,fill=BookPaper,line width=0.8pt,
    text width=30mm,minimum height=11mm,align=center,inner sep=2mm},
  train/.style={draw=BookTeal,fill=FigureModel!10,line width=0.8pt,
    text width=30mm,minimum height=11mm,align=center,inner sep=2mm},
  data/.style={draw=BookBlue,fill=FigureInput!10,line width=0.8pt,
    text width=21mm,minimum height=11mm,align=center,inner sep=2mm},
  add/.style={circle,draw=BookInk,fill=white,line width=0.8pt,
    minimum size=6mm,inner sep=0pt},
  flow/.style={draw=BookInk,line width=0.9pt,-{Latex},rounded corners=1mm},
  wire/.style={draw=BookInk,line width=0.9pt,rounded corners=1mm},
  heading/.style={anchor=west,font=\figurefont\small}
]
  \node[anchor=west,text=BookMuted] at (0,24)
    {共同条件：来源已有模型参数冻结；彩色模块：新增并训练的任务模型参数};
  \node[heading] at (0,13) {(a) 新建预测输出头：读取固定表示};
  \node[data] (head-input) at (13,0) {目标预测输入$\bm x$};
  \node[frozen] (backbone) at (61,0) {冻结底座\\$h_{\bm{\theta}_S}$};
  \node[train] (head) at (113,0) {训练新建头\\$\bm W_T,\bm b_T$};
  \node (head-output) at (157,0) {$\hat y$};
  \draw[flow] (head-input.east) -- (backbone.west);
  \draw[flow] (backbone.east) -- (head.west);
  \draw[flow] (head.east) -- (head-output.west);

  \node[heading] at (0,-22) {(b) 适配器：隐藏表示上的瓶颈残差};
  \node (adapter-input) at (10,-43) {$\bm h$};
  \coordinate (adapter-split) at (24,-43);
  \node[train] (down) at (52,-43) {训练降维\\再经非线性$\sigma$};
  \node[train] (up) at (103,-43) {训练升维\\$\bm W_{\mathrm{up}},\bm b_{\mathrm{up}}$};
  \node[add] (adapter-add) at (140,-43) {$+$};
  \node (adapter-output) at (158,-43) {$\bm h'$};
  \draw[wire] (adapter-input.east) -- (adapter-split);
  \draw[flow] (adapter-split) -- (down.west);
  \draw[flow] (down.east) -- (up.west);
  \draw[flow] (up.east) -- (adapter-add.west);
  \draw[flow] (adapter-split) -- (24,-31) -- (140,-31) -- (adapter-add.north);
  \draw[flow] (adapter-add.east) -- (adapter-output.west);

  \node[heading] at (0,-66) {(c) LoRA：线性映射上的并行低秩增量};
  \node (lora-input) at (10,-87) {$\bm h$};
  \coordinate (lora-split) at (24,-87);
  \node[frozen] (base-weight) at (78,-79) {冻结$\bm W_0$};
  \node[train] (low-rank) at (78,-99) {训练$s\bm B\bm A$};
  \node[add] (lora-add) at (140,-87) {$+$};
  \node (lora-output) at (158,-87) {$\bm h'$};
  \draw[wire] (lora-input.east) -- (lora-split);
  \draw[flow] (lora-split) |- (base-weight.west);
  \draw[flow] (lora-split) |- (low-rank.west);
  \draw[flow] (base-weight.east) -| (lora-add.north);
  \draw[flow] (low-rank.east) -| (lora-add.south);
  \draw[flow] (lora-add.east) -- (lora-output.west);

  \node[heading] at (0,-119) {(d) 软提示：预测输入序列中的连续任务条件};
  \node[train,text width=22mm] (prompt) at (15,-136) {训练提示\\$\bm P_T$};
  \node[data,text width=22mm] (tokens) at (15,-155) {词元嵌入\\$\bm E$};
  \node[data,text width=27mm] (concat) at (66,-145.5) {按序列位置拼接\\$[\bm P_T;\bm E]$};
  \node[frozen] (prompt-model) at (116,-145.5) {冻结模型\\词元长度$\ell+m$};
  \node (prompt-output) at (160,-145.5) {$\hat y$};
  \draw[flow] (prompt.east) -- (38,-136) |- (concat.west);
  \draw[flow] (tokens.east) -- (38,-155) |- (concat.west);
  \draw[flow] (concat.east) -- (prompt-model.west);
  \draw[flow] (prompt-model.east) -- (prompt-output.west);
\end{tikzpicture}

  \caption{来源权重冻结时，任务新增量进入计算的四个位置。实线为前向流，训练模块保存新增参数，冻结模块沿用来源权重，加号表示同维相加。适配器只画内部瓶颈残差，LoRA两支路共用同一输入；软提示与输入嵌入按位置拼接，连续提示不等同于现成文本示范。冻结模块仍参与前向，必要时也传递输入梯度。}
  \label{fig:ket-transfer-added-locations}
\end{figure*}
\FloatBarrier

\begin{example}[固定底座下的参数与计算核算]
  \label{ex:ket-transfer-parameter-counts}
  固定一个教学底座：12个Transformer层、隐藏宽度$d=768$、新任务5类。下面只核算指定配置的参数与注意力大小，不假设各方法风险相同。新线性头有$5(768+1)=3845$个参数。

  每层插两个$r=16$的适配器，则每个$2\cdot768\cdot16+16+768=25360$，24个共608640，加新头612485。若每层两个已有层归一化（LayerNorm）的缩放与偏置也可训练，再加$12\cdot2\cdot2\cdot768=36864$，总数649349。LayerNorm按每个位置的隐藏特征归一化，其缩放、偏置是原有量；因此后一配置同时包含已有部分更新。

  每层仅在查询和值两个$768\times768$映射上加$r=8$的LoRA，每矩阵12288个参数，共294912，加新头298757。每个未合并分支、每位置约增加12288次乘加；合并后免这条额外分支，但还需核算合并与切换。

  20个软提示增加15360个参数，复用现成生成输出时无需同一个分类头；若另加上述五类头，合计19205。长度128增至148时，单注意力头分数矩阵从16384增为21904个元素，多5520个，约33.69\%。这个比例只对应分数矩阵，不是完整前向时间或显存增幅。
\end{example}

若未来只允许训练小模块，还可在历史任务上准备它的适应接口。Meta-Adapters保留冻结底座$\bm\theta_S$、任务内可调普通适配器$\bm\phi$，并在其前后加入元适配器$\bm\omega$。历史内层改普通适配器，查询外层训练元适配器、普通适配器起点与逐层步长。为了让小模块经一步支持更新后降低查询损失，目标为
\begin{align}
  \bm\phi_k'
    &=\bm\phi_0-\bm\alpha\odot
       \nabla_{\bm\phi}L_{S_k}
         (\bm{\theta}_S,\bm\phi_0,\bm\omega), \nonumber\\
  \min_{\bm\omega,\bm\phi_0,\bm\alpha}\quad
  &\frac1M\sum_{k=1}^{M}
    L_{Q_k}(\bm{\theta}_S,\bm\phi_k',\bm\omega).
  \label{eq:ket-transfer-meta-adapters}
\end{align}
式中$\bm\alpha$把每层步长广播到该层适配器坐标，外层在规定的模块参数与步长范围内优化。目标端固定底座和元适配器，从$\bm\phi_0$出发只调整普通适配器；元适配器仍执行并参与输入梯度传播。因此迁移的既有小模块起点，也有帮助其适应的共享计算。共享$\bm\omega$与逐任务普通适配器须分开存储计数，实际配置的任务头也不能漏算\scite{ket-bansal2022metaadapters}

表~\ref{tab:ket-transfer-operation-cost}回到同一个新物种任务，分别列任务状态、一次适应与逐次预测。令$n$为支持数，$N$为类数，$d$为表示维数，$T$为更新次数，$b$为每次小批量大小，$F$为单输入底座前向费用。$B_{\mathcal A}$表示单样本反向费用加按样本分摊的参数更新费用，$\mathcal A$为可训练集合，因而$Tb(F+B_{\mathcal A})$只是一种主要项估算；批量硬件效率须另测。序列费用写成$F(\ell)$以显式保留长度，学习更新器与检索索引还须单列。

\begin{table*}[htbp]
  \centering
  \small
  \setlength{\tabcolsep}{4pt}
  \renewcommand{\arraystretch}{1.2}
  \begin{tabularx}{\linewidth}{@{}>{\raggedright\arraybackslash}p{24mm}
      >{\raggedright\arraybackslash}p{36mm}
      >{\raggedright\arraybackslash}X
      >{\raggedright\arraybackslash}X@{}}
    \toprule
    操作 & 保存的任务状态 & 一次性适应费用 & 逐次预测费用 \\
    \midrule
    语义参照 & 属性表或约$Nd$个向量参照分量 & 名称、属性或描述的整理与编码；不必优化模型 & 图像编码与$N$类比较；向量形式约$F+O(Nd)$ \\
    \addlinespace
    支持样本或原型 & 逐例约$nd$个表示分量，原型约$Nd$个；另存标签 & 提取$n$个带标签样本的特征约$nF$，原型汇总另需$O(nd)$ & 逐例约$F+O(nd)$，原型约$F+O(Nd)$ \\
    \addlinespace
    上下文条件化 & 许可的说明、示范与历史，或其摘要、键值缓存 & 检索与组织证据，计算摘要或固定前缀 & 依接口为$F(\ell)$；长上下文与生成另计 \\
    \addlinespace
    预测头拟合 & 已有头或新建头，约$N(d+1)$个参数 & 可缓存固定特征时约$nF+O(TbNd)$ & 底座前向与头部计算，约$F+O(Nd)$ \\
    \addlinespace
    全量微调 & 完整任务权重或等价差量，含必要的新头 & 约$Tb(F+B_{\mathcal A})$；$\mathcal A$含全部来源参数及新头 & 适应后完整模型前向 \\
    \addlinespace
    附加任务参数 & 适配器、低秩增量或软提示参数，以及底座版本 & 前向、相关反向与新参数更新；LoRA合并和切换另计 & 底座前向加模块或输入位置开销；LoRA合并后无额外分支 \\
    \bottomrule
  \end{tabularx}
  \caption{允许使用相同证据时，各操作的状态与主要计算项。方案可组合多行，也不必用尽证据；新头既是头拟合也是附加参数。支持比较采用固定独立编码，完整上下文编码需另加费用；头拟合假定可缓存特征，全量项采用单样本分摊估算，LoRA合并要求固定任务与兼容格式。来源编码器、语义模型和元训练另按式~\eqref{eq:ket-reuse-amortization}计入准备，不给出未经实测的性能排序。}
  \label{tab:ket-transfer-operation-cost}
\end{table*}

参数量、标签量、训练时间和预测费用回答四个不同问题。软提示参数少却须长路径反传，原型不做梯度更新却使用支持标签，全量训练的一次适应费用还可能被大量请求摊薄。判断整体效率，应固定允许使用的证据与目标标准，并把检索、提示组织、长度、更新、延迟与状态存储同列，不能以其中一项替代全部投入。

目标超出历史覆盖时，还要允许改变方案。来源只识别整图物种、目标却依赖微小叶斑，应先核对观测与表示是否保留细节，再结合增加标注和更新次数后的独立风险曲线定位困难。没有改善不唯一指向表示不足，也可能来自监督、优化或语义错误；核查后再补观测、扩大训练范围或回退普通目标训练。迁移收益应沿多个预算点与匹配对照比较，不只报告有利的一步结果。形成能力后继续服务的版本、记忆及费用，再交由第~\ref{sec:ket-update-maintenance}节持续维护。

\section{本章小结}
\label{sec:ket-transfer-summary}

回到新植物园，来源知识有用的前提是与目标存在可利用的联系。新物种需要类别含义与区分证据，病害需要观测和表示保留相关细节，定位需要对象级位置与匹配结构。拒识未知只提供目录不足时的退出，不会自动给新类命名、发现内部结构或建立分类能力。目标语义、许可信息和评价范围必须先明确，才能判断迁移减少了哪种错误。

准备与使用要相互对应：可比较表示由支持形成参照，图文联系让描述进入预测，起点、先验和更新器帮助参数估计，条件模型直接读取任务证据。元学习将接受任务证据后的结果纳入历史准备，不限定目标必须做几步梯度。目标端还应分别记录来源参数更新范围与新增参数；适配器、LoRA和软提示的位置、反向路径与部署状态不同，参数少不能成为统一的效率结论。

验证既要隔离支持与查询，也要隔离历史准备、选模和最终任务。固定证据、目标标准及预算，比较独立目标适应曲线与无迁移对照，才能说明来源帮助了学习；准备损失低或一次结果好都不足够。准备投入按实际使用摊销，失败与反复选择也计入账目。这样得到的是本次能力形成及成本的证据，旧能力保留和未来继续学习仍需各自的长期验证。
